% ROZDZIAŁ 2 — plik włączany przez main.tex
\chapter{Miary i ich podstawowe własności}
\label{ch:miary-wlasnosci}

\section{Definicja i pierwsze przykłady}

W rozdziale \ref{ch:przestrzenie-mierzalne} ustaliliśmy, które podzbiory
przestrzeni mają być mierzalne. Teraz przypiszemy im liczby z przedziału
$[0,\infty]$. Dopuszczenie wartości $\infty$ jest konieczne już dla
naturalnego pojęcia liczby elementów zbioru, a następnie dla miary
objętości na przestrzeniach nieograniczonych.

\begin{definicja}[Miara]
\label{def:measure}
Niech $(X,\mathcal F)$ będzie przestrzenią mierzalną.
\emph{Miarą} na $(X,\mathcal F)$ nazywamy funkcję
$\mu:\mathcal F\to[0,\infty]$ spełniającą:
\begin{enumerate}
\item $\mu(\varnothing)=0$;
\item dla każdego ciągu $(A_n)_{n\geq1}$ parami rozłącznych zbiorów
z $\mathcal F$ zachodzi \emph{przeliczalna addytywność}
\begin{equation}
\label{eq:countable-additivity}
\mu\left(\bigcup_{n=1}^{\infty}A_n\right)
=\sum_{n=1}^{\infty}\mu(A_n).
\end{equation}
\end{enumerate}
Trójkę $(X,\mathcal F,\mu)$ nazywamy \emph{przestrzenią z miarą}.
Wszystkie sumy liczb nieujemnych rozumiemy jako kres górny sum
częściowych, dopuszczając wartość $\infty$.
\end{definicja}

\begin{uwaga}
Nie jest wymagane, aby $\mu(X)$ było skończone. Aksjomat
$\mu(\varnothing)=0$ zapisujemy oddzielnie, ponieważ bez niego
funkcja stale równa $\infty$ spełniałaby równość
\eqref{eq:countable-additivity} dla każdej rodziny rozłącznej,
ale nie miałaby zamierzonych własności miary.
\end{uwaga}

\begin{przyklad}[Miara zerowa i miara Diraca]
\label{ex:dirac}
Na dowolnej przestrzeni mierzalnej $(X,\mathcal F)$ funkcja
$\mu(A)=0$ dla $A\in\mathcal F$ jest miarą zerową.
Dla ustalonego $x\in X$ definiujemy \emph{miarę Diraca} wzorem
\[
\delta_x(A)=\begin{cases}1,&x\in A,\\0,&x\notin A,\end{cases}
\qquad A\in\mathcal F.
\]
Jeśli $(A_n)$ są parami rozłączne, punkt $x$ należy co najwyżej
do jednego z tych zbiorów. Daje to przeliczalną addytywność.
Nie zakładamy przy tym, że singleton $\{x\}$ należy do $\mathcal F$.
\end{przyklad}

\begin{przyklad}[Miara licząca]
\label{ex:counting}
Na $(X,\mathcal P(X))$ definiujemy
\[
\#(A)=\begin{cases}
|A|,&A\text{ jest skończony},\\
\infty,&A\text{ jest nieskończony}.
\end{cases}
\]
Dla parami rozłącznych zbiorów $A_n$ liczność ich sumy jest sumą
liczności w podanym sensie. Jeśli $\bigcup_n A_n$ jest skończony,
tylko skończenie wiele zbiorów $A_n$ jest niepustych i każdy
jest skończony, więc równość addytywności wynika ze
zwykłego zliczania elementów. Jeśli $\bigcup_n A_n$
jest nieskończony, to albo pewien $A_n$ jest nieskończony,
albo nieskończenie wiele zbiorów $A_n$ jest niepustych.
W obu przypadkach $\sum_n\#(A_n)=\infty=\#(\bigcup_nA_n)$.
Jest to \emph{miara licząca}.
\end{przyklad}

\begin{przyklad}[Miara na $\sigma$-algebrze przeliczalno-koprzeliczalnej]
\label{ex:cocountable-measure}
Niech $X$ będzie nieprzeliczalnym zbiorem i niech
$\mathcal C_X$ będzie $\sigma$-algebrą z przykładu
\ref{ex:countable-cocountable}. Określmy
\[
\mu(A)=\begin{cases}
0,&A\text{ jest przeliczalny},\\
1,&X\setminus A\text{ jest przeliczalny}.
\end{cases}
\qquad A\in\mathcal C_X.
\]
Definicja jest jednoznaczna, ponieważ nieprzeliczalny zbiór
$X$ nie może być jednocześnie przeliczalny i mieć przeliczalne
dopełnienie. Dwie rozłączne części $X$ nie mogą obie mieć
przeliczalnego dopełnienia. Dla przeliczalnej rodziny
rozłącznych zbiorów z $\mathcal C_X$ albo wszystkie mają miarę
zero, a ich suma jest przeliczalna, albo dokładnie jeden ma
miarę jeden, a suma ma miarę jeden. Jest to zatem miara.
Ten przykład pokaże w sekcji \ref{sec:atoms}, że atomy
nie muszą być singletonami.
\end{przyklad}

\begin{przyklad}[Ważone miary dyskretne]
\label{ex:weighted-discrete}
Niech $X$ będzie zbiorem, $\mathcal F=\mathcal P(X)$,
$(x_n)_{n\geq1}$ ciągiem punktów $X$, a $(w_n)_{n\geq1}$
ciągiem liczb z $[0,\infty)$. Wtedy
\[
\mu(A)=\sum_{n=1}^{\infty}w_n\chi_A(x_n),
\qquad A\subseteq X,
\]
jest miarą. Dla parami rozłącznych $A_k$ mamy
$\chi_{\bigcup_k A_k}(x_n)=\sum_k\chi_{A_k}(x_n)$.
Możemy zamienić kolejność sumowania, ponieważ dla
$a_{n,k}=w_n\chi_{A_k}(x_n)\geq0$ obie sumy iterowane
$\sum_n\sum_k a_{n,k}$ i $\sum_k\sum_n a_{n,k}$
są równe kresowi górnemu sum
$\sum_{n\in F}\sum_{k\in G}a_{n,k}$
po wszystkich skończonych zbiorach indeksów $F,G\subseteq\mathbb N$.
Otrzymujemy więc przeliczalną addytywność, także gdy
któraś suma jest nieskończona. Nie wymagamy, aby punkty
$x_n$ były różne ani by $\sum_nw_n$ było skończone.
\end{przyklad}

\section{Monotoniczność i przeliczalna podaddytywność}

\begin{twierdzenie}[Podstawowe własności miary]
\label{thm:basic-measure}
Niech $(X,\mathcal F,\mu)$ będzie przestrzenią z miarą.
Wówczas:
\begin{enumerate}
\item dla dowolnych $A,B\in\mathcal F$ z $A\subseteq B$ zachodzi
$\mu(A)\leq\mu(B)$;
\item dla dowolnych $A,B\in\mathcal F$ zachodzi
$\mu(A\cup B)\leq\mu(A)+\mu(B)$;
\item dla dowolnych $A_n\in\mathcal F$, $n\geq1$, zachodzi
\begin{equation}
\label{eq:subadditivity}
\mu\left(\bigcup_{n=1}^{\infty}A_n\right)
\leq\sum_{n=1}^{\infty}\mu(A_n);
\end{equation}
\item jeśli $A,B\in\mathcal F$, $A\subseteq B$ i $\mu(A)<\infty$,
to $\mu(B\setminus A)=\mu(B)-\mu(A)$.
\end{enumerate}
\end{twierdzenie}
\begin{proof}
Zbiory $A$ i $B\setminus A$ są rozłączne, a ich suma jest równa $B$.
Przeliczalna addytywność, zastosowana do rodziny uzupełnionej
zbiorami pustymi, daje
$\mu(B)=\mu(A)+\mu(B\setminus A)$. Wynika stąd monotoniczność,
a jeżeli $\mu(A)<\infty$, również wzór na różnicę.

Ponadto $A\cup B=A\,\dot\cup\,(B\setminus A)$, więc
$\mu(A\cup B)\leq\mu(A)+\mu(B)$.
Dla dowolnego ciągu $(A_n)$ definiujemy
$B_1=A_1$ oraz
$B_n=A_n\setminus\bigcup_{k<n}A_k$ dla $n\geq2$.
Zbiory $B_n$ są parami rozłączne, $B_n\subseteq A_n$
i $\bigcup_nB_n=\bigcup_nA_n$. Stąd
$\mu(\bigcup_nA_n)=\sum_n\mu(B_n)\leq\sum_n\mu(A_n)$.
\end{proof}

\begin{uwaga}[Arytmetyka z nieskończonością]
Warunku $\mu(A)<\infty$ w ostatnim punkcie twierdzenia
\ref{thm:basic-measure} nie wolno usuwać. Jeśli
$\mu(A)=\mu(B)=\infty$, wyrażenie $\infty-\infty$
nie ma określonej wartości. Prawidłowym wzorem pozostaje
$\mu(B)=\mu(A)+\mu(B\setminus A)$, który nie wymaga odejmowania.
\end{uwaga}

\begin{wniosek}[Przeliczalne sumy zbiorów miary zero]
\label{cor:null-union}
Jeżeli $(X,\mathcal F,\mu)$ jest przestrzenią z miarą oraz
$N_n\in\mathcal F$ i $\mu(N_n)=0$ dla każdego $n\in\mathbb N$,
to $\mu(\bigcup_nN_n)=0$.
\end{wniosek}
\begin{proof}
Z podaddytywności mamy
$0\leq\mu(\bigcup_nN_n)\leq\sum_n\mu(N_n)=0$.
\end{proof}

\section{Ciągłość miary względem granic zbiorów}

\begin{twierdzenie}[Ciągłość z dołu i z góry]
\label{thm:measure-continuity}
Niech $(X,\mathcal F,\mu)$ będzie przestrzenią z miarą.
\begin{enumerate}
\item Jeśli $A_n\in\mathcal F$ oraz $A_n\uparrow A$, to
$A\in\mathcal F$ i
\begin{equation}
\label{eq:continuity-below}
\mu(A)=\lim_{n\to\infty}\mu(A_n).
\end{equation}
\item Jeśli $B_n\in\mathcal F$, $B_n\downarrow B$ oraz
$\mu(B_1)<\infty$, to $B\in\mathcal F$ i
\begin{equation}
\label{eq:continuity-above}
\mu(B)=\lim_{n\to\infty}\mu(B_n).
\end{equation}
\end{enumerate}
\end{twierdzenie}
\begin{proof}
Dla ciągu rosnącego przyjmujemy $E_1=A_1$ i
$E_n=A_n\setminus A_{n-1}$ dla $n\geq2$.
Zbiory $E_n$ są parami rozłączne oraz
$A=\bigcup_nE_n$ i $A_k=\bigcup_{n=1}^kE_n$.
Stąd
$\mu(A)=\sum_n\mu(E_n)=\lim_{k\to\infty}\sum_{n=1}^k\mu(E_n)
=\lim_k\mu(A_k)$.

Dla ciągu malejącego zbiory $C_n=B_1\setminus B_n$ rosną
do $B_1\setminus B$. Z pierwszej części mamy
$\mu(C_n)\to\mu(B_1\setminus B)$. Ponieważ
$\mu(B_1)<\infty$, wolno zastosować wzór na różnicę:
$\mu(C_n)=\mu(B_1)-\mu(B_n)$ oraz
$\mu(B_1\setminus B)=\mu(B_1)-\mu(B)$.
Po odjęciu od $\mu(B_1)$ otrzymujemy żądaną równość.
\end{proof}

\begin{uwaga}[Wystarczy skończoność miary jednego wyrazu]
W drugim punkcie twierdzenia~\ref{thm:measure-continuity} można
zastąpić założenie $\mu(B_1)<\infty$ słabszym
$\mu(B_m)<\infty$ dla pewnego $m\ge1$: wystarczy
zastosować twierdzenie do ogona ciągu
$(B_{m+k-1})_{k\ge1}$, którego przekrój jest nadal równy $B$.
Nie wystarcza natomiast sama skończoność $\mu(B)$.
\end{uwaga}

\begin{przyklad}[Dlaczego założenie skończoności jest niezbędne]
\label{ex:continuity-counterexample}
Na $X=\mathbb N$ z miarą liczącą zbiory
$B_n=\{n,n+1,n+2,\ldots\}$ tworzą ciąg malejący o pustym
przekroju. Każdy $B_n$ ma miarę $\infty$, ale
$\mu(\bigcap_nB_n)=\mu(\varnothing)=0$.
Tak więc ciągłość z góry nie zachodzi bez odpowiedniej
hipotezy skończoności.
\end{przyklad}

\begin{definicja}[Granica górna i dolna zbiorów]
\label{def:limsup-sets}
Dla ciągu $(A_n)_{n\geq1}$ podzbiorów $X$ definiujemy
\[
\liminf_{n\to\infty}A_n
=\bigcup_{m=1}^{\infty}\bigcap_{n\geq m}A_n,
\qquad
\limsup_{n\to\infty}A_n
=\bigcap_{m=1}^{\infty}\bigcup_{n\geq m}A_n.
\]
Pierwszy zbiór składa się z punktów należących do wszystkich
$A_n$ od pewnego miejsca, drugi --- z punktów należących do
nieskończenie wielu $A_n$. Jeśli oba zbiory są równe,
ich wspólną wartość oznaczamy przez $\lim_nA_n$.
\end{definicja}

\begin{twierdzenie}[Nierówności dla granic zbiorów]
\label{thm:fatou-for-sets}
Niech $(X,\mathcal F,\mu)$ będzie przestrzenią z miarą
oraz $A_n\in\mathcal F$ dla wszystkich $n\ge1$.
Wówczas
\[
 \mu(\liminf_{n\to\infty}A_n)
 \le \liminf_{n\to\infty}\mu(A_n).
\]
Jeżeli dodatkowo $\mu(X)<\infty$, to
\[
 \limsup_{n\to\infty}\mu(A_n)
 \le \mu(\limsup_{n\to\infty}A_n).
\]
\end{twierdzenie}
\begin{proof}
Niech $C_m=\bigcap_{n\ge m}A_n$.
Zbiory $C_m$ tworzą ciąg rosnący,
$\bigcup_m C_m=\liminf_n A_n$ oraz
$\mu(C_m)\le\inf_{n\ge m}\mu(A_n)$.
Ciągłość miary z dołu daje
$\mu(\liminf_n A_n)=\lim_m\mu(C_m)
\le\lim_m\inf_{n\ge m}\mu(A_n)$.
Dla drugiej nierówności kładziemy
$D_m=\bigcup_{n\ge m}A_n$.
Zbiory $D_m$ maleją do $\limsup_n A_n$,
a $\mu(D_m)\ge\sup_{n\ge m}\mu(A_n)$.
Ponieważ $\mu(D_1)\le\mu(X)<\infty$,
ciągłość miary z góry pozwala przejść do granicy:
$\mu(\limsup_n A_n)=\lim_m\mu(D_m)
\ge\lim_m\sup_{n\ge m}\mu(A_n)$.
\end{proof}
\begin{uwaga}
Założenie skończoności jest istotne w drugiej nierówności.
Dla miary liczącej na $\mathbb N$ i
$A_n=\{n,n+1,\ldots\}$ mamy
$\limsup_n\mu(A_n)=\infty$, lecz
$\limsup_n A_n=\varnothing$.
\end{uwaga}

\begin{twierdzenie}[Lemat Borela--Cantellego, część pierwsza]
\label{thm:borel-cantelli-1}
Niech $(X,\mathcal F,\mu)$ będzie przestrzenią z miarą
oraz $(A_n)_{n\geq1}$ ciągiem zbiorów z $\mathcal F$.
Jeżeli $\sum_{n=1}^{\infty}\mu(A_n)<\infty$, to
$\mu(\limsup_{n\to\infty}A_n)=0$.
\end{twierdzenie}
\begin{proof}
Zbiór $\limsup_nA_n$ jest mierzalny i dla każdego
$m\in\mathbb N$ zawiera się w $\bigcup_{n\geq m}A_n$.
Z monotoniczności i podaddytywności wynika
\[
0\leq\mu(\limsup_nA_n)
\leq\sum_{n\geq m}\mu(A_n).
\]
Ponieważ ogólny szereg jest zbieżny, jego ogon dąży do zera,
gdy $m\to\infty$. To kończy dowód. Nie użyliśmy ani
skończoności $\mu(X)$, ani niezależności zbiorów.
\end{proof}

\begin{przyklad}[Interpretacja]
Jeśli $\mu(A_n)\leq2^{-n}$ dla wszystkich $n\geq1$, to
$\mu(\limsup_nA_n)=0$. Prawie każdy punkt (w sensie, który
zaraz zdefiniujemy) może należeć tylko do skończenie wielu $A_n$.
W rachunku prawdopodobieństwa po wprowadzeniu niezależności
uzyskuje się drugą część lematu Borela--Cantellego, z innymi
założeniami i przeciwnym wnioskiem.
\end{przyklad}

\section{Skończoność, \texorpdfstring{$\sigma$}{sigma}-skończoność i półskończoność}

\begin{definicja}[Klasy miar]
\label{def:measure-finiteness}
Niech $(X,\mathcal F,\mu)$ będzie przestrzenią z miarą.
\begin{enumerate}
\item Miara $\mu$ jest \emph{skończona}, gdy $\mu(X)<\infty$.
\item Miara $\mu$ jest \emph{probabilistyczna}, gdy $\mu(X)=1$.
\item Miara $\mu$ jest \emph{$\sigma$-skończona}, gdy istnieją
$E_n\in\mathcal F$, $n\geq1$, takie że
$X=\bigcup_nE_n$ i $\mu(E_n)<\infty$ dla wszystkich $n$.
\item Miara $\mu$ jest \emph{półskończona}, gdy dla każdego
$A\in\mathcal F$ z $\mu(A)>0$ istnieje $B\in\mathcal F$
takie, że $B\subseteq A$ oraz $0<\mu(B)<\infty$.
\end{enumerate}
Jeśli $(X,\tau)$ jest przestrzenią topologiczną, a
$\mathcal F\supseteq\mathcal B(X)$, nazywamy $\mu$
\emph{lokalnie skończoną}, gdy każdy $x\in X$ posiada
otwarte otoczenie $U$ o $\mu(U)<\infty$.
\end{definicja}

\begin{twierdzenie}[Implikacje]
\label{thm:finiteness-implications}
Każda miara skończona na przestrzeni mierzalnej jest
$\sigma$-skończona, a każda miara $\sigma$-skończona
jest półskończona. Implikacje odwrotne są na ogół fałszywe.
\end{twierdzenie}
\begin{proof}
Jeśli $\mu(X)<\infty$, można w definicji $\sigma$-skończoności
przyjąć $E_n=X$. Niech teraz $X=\bigcup_nE_n$ z
$\mu(E_n)<\infty$. Dla $A\in\mathcal F$ o $\mu(A)>0$
nie wszystkie przekroje $A\cap E_n$ mogą mieć miarę zero:
w przeciwnym razie podaddytywność dałaby
$\mu(A)\leq\sum_n\mu(A\cap E_n)=0$.
Istnieje więc $n$ z $0<\mu(A\cap E_n)<\infty$.

Miara licząca na $\mathbb N$ jest $\sigma$-skończona,
ale nieskończona. Miara licząca na nieprzeliczalnym
zbiorze $X$ jest półskończona (każdy niepusty zbiór
zawiera singleton miary jeden), ale nie jest
$\sigma$-skończona, gdyż zbiory o skończonej mierze
są skończone, a przeliczalna suma zbiorów skończonych
jest przeliczalna.
\end{proof}

\begin{przyklad}[Miara, która nie jest półskończona]
Niech $X$ będzie niepustym zbiorem z $\mathcal F=\mathcal P(X)$
i niech $\mu(\varnothing)=0$ oraz $\mu(A)=\infty$
dla $A\ne\varnothing$. Jest to miara: dla rodziny
rozłącznej z niepustą sumą co najmniej jeden składnik
ma miarę $\infty$. Nie ma jednak żadnego podzbioru
o mierze dodatniej i skończonej.
\end{przyklad}

\begin{uwaga}[Lokalna skończoność a $\sigma$-skończoność]
Na przestrzeni topologicznej mającej przeliczalną bazę
lokalnie skończona miara borelowska jest $\sigma$-skończona.
Istotnie, otwarte otoczenia o skończonej mierze tworzą
pokrycie; własność Lindelöfa (wynikająca z istnienia
przeliczalnej bazy) pozwala wybrać pokrycie przeliczalne.
Bez założenia topologicznego takiej implikacji nie ma:
miara licząca na nieprzeliczalnej przestrzeni dyskretnej
jest lokalnie skończona, lecz nie jest $\sigma$-skończona.
\end{uwaga}

\section{Zbiory miary zero, własności prawie wszędzie i zupełność}

\begin{definicja}[Własność prawie wszędzie]
\label{def:almost-everywhere}
Niech $(X,\mathcal F,\mu)$ będzie przestrzenią z miarą.
Mówimy, że własność $P(x)$ zachodzi \emph{$\mu$-prawie wszędzie}
(skrót: $\mu$-p.w.), jeśli istnieje $N\in\mathcal F$
o $\mu(N)=0$ takie, że $P(x)$ zachodzi dla każdego
$x\in X\setminus N$. Nie wymagamy mierzalności zbioru
$\{x\in X:P(x)\text{ nie zachodzi}\}$.
Jeśli jest on mierzalny, jest to równoważne temu,
że ma miarę zero.
\end{definicja}

\begin{uwaga}
Dzięki wnioskowi \ref{cor:null-union} przeliczalnie wiele
własności zachodzących $\mu$-p.w. zachodzi jednocześnie
$\mu$-p.w. Bez przeliczalności nie można wyciągać tego
wniosku: przy miarze z przykładu \ref{ex:cocountable-measure}
każdy singleton ma miarę zero, lecz suma wszystkich singletonów
jest równa $X$ i ma miarę jeden.
\end{uwaga}

\begin{definicja}[Miara zupełna]
\label{def:complete-measure}
Miara $\mu$ na $(X,\mathcal F)$ jest \emph{zupełna}, jeżeli
dla każdego $N\in\mathcal F$ spełniającego $\mu(N)=0$
każdy podzbiór $A\subseteq N$ należy do $\mathcal F$.
Z monotoniczności wynika wówczas $\mu(A)=0$.
\end{definicja}

\begin{przyklad}[Miara niezupełna na trzech punktach]
\label{ex:incomplete-measure}
Niech $X=\{a,b,c\}$ oraz
$\mathcal F=\{\varnothing,\{a,b\},\{c\},X\}$.
Określmy $\mu(\{a,b\})=0$ i $\mu(\{c\})=1$,
a pozostałe wartości wyznaczmy przez addytywność:
$\mu(\varnothing)=0$, $\mu(X)=1$.
To dobrze określona miara na $\mathcal F$.
Zbiór $N=\{a,b\}$ ma miarę zero, ale jego podzbiór
$\{a\}$ nie jest mierzalny. Zatem $\mu$ nie jest zupełna.
\end{przyklad}

\begin{twierdzenie}[Uzupełnienie przestrzeni z miarą]
\label{thm:completion}
Niech $(X,\mathcal F,\mu)$ będzie przestrzenią z miarą.
Zdefiniujmy
\[
\overline{\mathcal F}
=\{E\subseteq X:\text{istnieją }A,B\in\mathcal F
\text{ takie, że }A\subseteq E\subseteq B
\text{ oraz }\mu(B\setminus A)=0\}.
\]
Wówczas $\overline{\mathcal F}$ jest $\sigma$-algebrą zawierającą
$\mathcal F$, a wzór
\[
\overline\mu(E)=\mu(A),\qquad
A\subseteq E\subseteq B,\quad A,B\in\mathcal F,
\quad\mu(B\setminus A)=0,
\]
definiuje na niej miarę zupełną. Miara $\overline\mu$
przedłuża $\mu$, a $\overline{\mathcal F}$ jest najmniejszą
$\sigma$-algebrą zawierającą $\mathcal F$ i wszystkie
podzbiory zbiorów z $\mathcal F$ o mierze zero.
\end{twierdzenie}
\begin{proof}
Najpierw sprawdzamy jednoznaczność definicji $\overline\mu$.
Jeśli $A_i\subseteq E\subseteq B_i$ i
$\mu(B_i\setminus A_i)=0$ dla $i=1,2$, to
$A_1\setminus A_2\subseteq B_2\setminus A_2$.
Z monotoniczności mamy $\mu(A_1\setminus A_2)=0$,
a stąd $\mu(A_1)\leq\mu(A_2)$, nawet jeśli
jedna z miar jest nieskończona. Zamiana indeksów
daje $\mu(A_1)=\mu(A_2)$. Ponadto
$\mu(A_i)=\mu(B_i)$, gdyż
$B_i=A_i\,\dot\cup\,(B_i\setminus A_i)$.

Jeśli $E\in\overline{\mathcal F}$ i
$A\subseteq E\subseteq B$ są jak w tezie, to
$X\setminus B\subseteq X\setminus E\subseteq X\setminus A$,
a różnica zbiorów skrajnych wynosi $B\setminus A$
i ma miarę zero. To dowodzi zamknięcia na dopełnienia.
Jeśli $E_n\in\overline{\mathcal F}$ i
$A_n\subseteq E_n\subseteq B_n$ z
$\mu(B_n\setminus A_n)=0$, to
\[
\bigcup_nA_n\subseteq\bigcup_nE_n\subseteq\bigcup_nB_n,
\qquad
\left(\bigcup_nB_n\right)\setminus\left(\bigcup_nA_n\right)
\subseteq\bigcup_n(B_n\setminus A_n).
\]
Ostatni zbiór ma miarę zero na mocy wniosku
\ref{cor:null-union}, więc przeliczalna suma należy do
$\overline{\mathcal F}$. Zawieranie $\mathcal F$ wynika
z wyboru $A=B=E\in\mathcal F$.

Dla parami rozłącznych $E_n\in\overline{\mathcal F}$
wybieramy $A_n\subseteq E_n$ jak wyżej. Zbiory $A_n$
są parami rozłączne i
$\overline\mu(\bigcup_nE_n)=\mu(\bigcup_nA_n)
=\sum_n\mu(A_n)=\sum_n\overline\mu(E_n)$.
Zatem $\overline\mu$ jest miarą i przedłuża $\mu$.

Aby wykazać zupełność, niech $E\in\overline{\mathcal F}$
i $\overline\mu(E)=0$. Wybierzmy $A,B\in\mathcal F$
takie, że $A\subseteq E\subseteq B$ i
$\mu(B\setminus A)=0$. Mamy $\mu(B)=\mu(A)=0$.
Dla dowolnego $T\subseteq E$ zachodzi
$\varnothing\subseteq T\subseteq B$, przy $\mu(B)=0$,
więc $T\in\overline{\mathcal F}$.

Wreszcie, jeśli $\mathcal G$ jest $\sigma$-algebrą zawierającą
$\mathcal F$ i wszystkie podzbiory jej zbiorów miary zero,
to dla $E\in\overline{\mathcal F}$ jak wyżej zbiór
$E\setminus A$ jest podzbiorem mierzalnego zbioru
$B\setminus A$ miary zero. Zatem $E\setminus A\in\mathcal G$
i $E=A\cup(E\setminus A)\in\mathcal G$. Dowodzi to minimalności.
\end{proof}

\begin{uwaga}[Jednoznaczność uzupełnienia]
Konstrukcja z twierdzenia \ref{thm:completion} nie polega
na arbitralnym wyborze wartości dla nowych podzbiorów.
Zbiory nowo dołączone różnią się od starych zbiorów
mierzalnych jedynie o części zbiorów miary zero.
Rozróżnienie między miarą borelowską a jej uzupełnieniem
będzie ważne przy konstrukcji miary Lebesgue'a.
\end{uwaga}

\section{Operacje na miarach}

\begin{twierdzenie}[Sumy i mnożenie przez stałą]
\label{thm:measure-operations}
Niech $(X,\mathcal F)$ będzie przestrzenią mierzalną,
$(\mu_n)_{n\geq1}$ ciągiem miar na $\mathcal F$,
a $(c_n)_{n\geq1}$ ciągiem liczb nieujemnych i skończonych.
Wówczas
$\mu(A)=\sum_{n=1}^{\infty}c_n\mu_n(A)$ dla $A\in\mathcal F$
jest miarą, przy konwencji, że $0\mu_n$ oznacza miarę
zerową również wtedy, gdy $\mu_n(A)=\infty$.
\end{twierdzenie}
\begin{proof}
Wszystkie składniki są nieujemne i wartość na zbiorze pustym
wynosi zero. Dla parami rozłącznych $A_k$ otrzymujemy
\[
\mu\left(\bigcup_kA_k\right)
=\sum_n c_n\sum_k\mu_n(A_k)
=\sum_k\sum_n c_n\mu_n(A_k)
=\sum_k\mu(A_k).
\]
Wolno przestawić kolejność sumowania, ponieważ obie strony
środkowej równości są kresami górnymi tych samych sum
skończonych wyrazów nieujemnych.
\end{proof}

\begin{definicja}[Miara śladowa i ograniczenie]
\label{def:restriction}
Niech $(X,\mathcal F,\mu)$ będzie przestrzenią z miarą
i $E\in\mathcal F$. Na przestrzeni mierzalnej
$(E,\mathcal F|_E)$ definiujemy \emph{miarę śladową}
$\mu|_E$ wzorem
$(\mu|_E)(B)=\mu(B)$ dla $B\in\mathcal F|_E$.
Ponieważ $E\in\mathcal F$, każdy zbiór śladowy
$B=A\cap E$ jest elementem $\mathcal F$.
Można też określić na $(X,\mathcal F)$ miarę
$\mu_E(A)=\mu(A\cap E)$; obie konstrukcje mają
różne dziedziny i nie należy ich utożsamiać.
\end{definicja}

\begin{proof}[Sprawdzenie aksjomatów]
Miary obu typów przypisują zbiorowi pustemu wartość zero.
Jeśli $B_n\subseteq E$ są parami rozłączne, to ich
suma jest rozłączna również w $X$, a przeliczalna
addytywność $\mu$ daje addytywność $\mu|_E$.
Dla $\mu_E$ korzystamy z tożsamości
$(\bigcup_nA_n)\cap E=\bigcup_n(A_n\cap E)$.
\end{proof}

\begin{definicja}[Miara obrazowa]
\label{def:pushforward}
Niech $(X,\mathcal F,\mu)$ będzie przestrzenią z miarą,
$(Y,\mathcal G)$ przestrzenią mierzalną, a
$T:X\to Y$ odwzorowaniem $(\mathcal F,\mathcal G)$-mierzalnym.
\emph{Miarą obrazową} $T_\#\mu$ na $(Y,\mathcal G)$ nazywamy
miarę określoną wzorem
\[
(T_\#\mu)(B)=\mu(T^{-1}(B)),\qquad B\in\mathcal G.
\]
\end{definicja}

\begin{proof}[Sprawdzenie aksjomatów]
Z mierzalności $T$ wszystkie obrazy odwrotne występujące
we wzorze są mierzalne. Obraz odwrotny zbioru pustego
jest pusty. Dla parami rozłącznych $B_n$ zbiory
$T^{-1}(B_n)$ są parami rozłączne i
$T^{-1}(\bigcup_nB_n)=\bigcup_nT^{-1}(B_n)$.
Równość przeliczalnej addytywności jest zatem konsekwencją
tego samego aksjomatu dla $\mu$.
\end{proof}

\begin{przyklad}[Obraz miary Diraca]
Jeżeli $x\in X$ i $T:X\to Y$ jest mierzalne, to
$T_\#\delta_x=\delta_{T(x)}$.
W szczególności obraz miary dyskretnej
$\sum_nw_n\delta_{x_n}$ jest równy
$\sum_nw_n\delta_{T(x_n)}$; punkty o tym samym obrazie
mogą powodować sumowanie wag.
\end{przyklad}

\begin{przyklad}[Normalizacja]
Jeżeli miara $\mu$ na $(X,\mathcal F)$ spełnia
$0<\mu(X)<\infty$, to $\mathbb P(A)=\mu(A)/\mu(X)$
jest miarą probabilistyczną. Nie można tego zrobić
w ten sposób dla miary o całkowitej masie nieskończonej.
\end{przyklad}

\section{Jednoznaczność miar na układach generujących}

\begin{twierdzenie}[Jednoznaczność miar skończonych]
\label{thm:uniqueness-finite}
Niech $X$ będzie zbiorem, $\mathcal P$ układem $\pi$
podzbiorów $X$, a $\mathcal F=\sigma(\mathcal P)$.
Jeżeli $\mu,\nu$ są miarami skończonymi na $(X,\mathcal F)$
spełniającymi $\mu(X)=\nu(X)$ i
$\mu(P)=\nu(P)$ dla każdego $P\in\mathcal P$,
to $\mu=\nu$ na całej $\mathcal F$.
\end{twierdzenie}
\begin{proof}
Niech $\mathcal D=\{A\in\mathcal F:\mu(A)=\nu(A)\}$.
Z założenia $X\in\mathcal D$ i $\mathcal P\subseteq\mathcal D$.
Jeżeli $A,B\in\mathcal D$ i $A\subseteq B$, to
$\mu(B\setminus A)=\mu(B)-\mu(A)
=\nu(B)-\nu(A)=\nu(B\setminus A)$.
Odejmowanie jest dopuszczalne, gdyż obie miary są skończone.
Dla parami rozłącznych $A_n\in\mathcal D$ mamy
$\mu(\bigcup_nA_n)=\sum_n\mu(A_n)
=\sum_n\nu(A_n)=\nu(\bigcup_nA_n)$.
Zatem $\mathcal D$ jest układem Dynkina zawierającym
$\mathcal P$. Twierdzenie \ref{thm:pi-lambda} daje
$\mathcal F\subseteq\mathcal D$.
\end{proof}

\begin{twierdzenie}[Jednoznaczność przy wspólnym wyczerpaniu]
\label{thm:uniqueness-sigma-finite}
Niech $X$ będzie zbiorem, $\mathcal P$ układem $\pi$ na $X$,
$\mathcal F=\sigma(\mathcal P)$, a $\mu,\nu$ miarami
na $(X,\mathcal F)$ zgodnymi na $\mathcal P$.
Załóżmy, że istnieją zbiory $E_n\in\mathcal P$ takie, że
$E_n\subseteq E_{n+1}$, $\bigcup_nE_n=X$ oraz
$\mu(E_n)=\nu(E_n)<\infty$ dla każdego $n$.
Wtedy $\mu=\nu$ na $\mathcal F$.
\end{twierdzenie}
\begin{proof}
Dla każdego $n$ definiujemy miary skończone na $\mathcal F$:
$\mu_n(A)=\mu(A\cap E_n)$ i
$\nu_n(A)=\nu(A\cap E_n)$.
Ponieważ $E_n\in\mathcal P$ i $\mathcal P$ jest układem $\pi$,
dla każdego $P\in\mathcal P$ zachodzi
$P\cap E_n\in\mathcal P$, a więc
$\mu_n(P)=\nu_n(P)$. Ponadto
$\mu_n(X)=\nu_n(X)<\infty$. Z twierdzenia
\ref{thm:uniqueness-finite} wynika $\mu_n=\nu_n$.
Dla dowolnego $A\in\mathcal F$ zbiory
$A\cap E_n$ rosną do $A$. Ciągłość miary z dołu
daje
$\mu(A)=\lim_n\mu(A\cap E_n)
=\lim_n\nu(A\cap E_n)=\nu(A)$.
\end{proof}

\begin{przyklad}[Miara określona wartościami dystrybuanty]
\label{ex:unique-cdf}
Niech $\mu$ i $\nu$ będą miarami probabilistycznymi
na $(\mathbb R,\mathcal B(\mathbb R))$.
Jeśli dla każdego $t\in\mathbb R$ zachodzi
$\mu((-\infty,t])=\nu((-\infty,t])$, to $\mu=\nu$.
Wystarczy bowiem zastosować twierdzenie
\ref{thm:uniqueness-finite} do układu $\pi$
$\{(-\infty,t]:t\in\mathbb R\}$.
Układ ten generuje $\mathcal B(\mathbb R)$,
a równość całkowitych mas wynika z probabilistyczności.
To wyjaśnia, dlaczego dystrybuanta jednoznacznie
wyznacza rozkład zmiennej losowej.
\end{przyklad}

\begin{przyklad}[Dlaczego wartości $\infty$ nie wystarczają]
\label{ex:uniqueness-fails}
Na $X=\mathbb R$ wyposażonym w $\mathcal B(\mathbb R)$
niech $\mu$ będzie miarą liczącą, a $\nu=2\mu$.
Obie miary przypisują wartość $\infty$ każdemu niepustemu
otwartemu przedziałowi o końcach wymiernych oraz wartość
zero zbiorowi pustemu. Zgadzają się zatem na układzie
$\pi$ tych przedziałów (razem ze zbiorem pustym),
który generuje $\mathcal B(\mathbb R)$.
Jednak $\mu(\{0\})=1$, a $\nu(\{0\})=2$.
Nie ma tu wspólnego wyczerpania zbiorami z układu $\pi$
o skończonych miarach. Przykład dowodzi, że zgodność
na generatorach bez odpowiednich warunków skończoności
nie wyznacza miary jednoznacznie.
\end{przyklad}

\section{Atomy miar}
\label{sec:atoms}

\begin{definicja}[Atom i miara bezatomowa]
\label{def:atom}
Niech $(X,\mathcal F,\mu)$ będzie przestrzenią z miarą.
Zbiór $A\in\mathcal F$ nazywamy \emph{atomem miary $\mu$},
jeżeli $\mu(A)>0$ oraz dla każdego $B\in\mathcal F$
zawartego w $A$ zachodzi $\mu(B)=0$ lub
$\mu(A\setminus B)=0$.
Miarę $\mu$ nazywamy \emph{bezatomową}, jeśli nie ma atomów.
\end{definicja}

\begin{przyklad}[Atomy miary dyskretnej]
Dla miary Diraca $\delta_x$ każdy mierzalny zbiór
zawierający $x$ jest atomem; wszystkie takie zbiory
są równoważne z dokładnością do zbiorów miary zero.
Dla miary liczącej na $\mathcal P(X)$ atomami są
dokładnie singletony: zbiór liczący co najmniej dwa punkty
można rozłożyć na dwie niepuste mierzalne części
o dodatniej mierze.
\end{przyklad}

\begin{przyklad}[Atom, który nie jest singletonem]
W przestrzeni z przykładu \ref{ex:cocountable-measure}
cały zbiór $X$ jest atomem: jego mierzalne podzbiory
mają miarę zero albo ich dopełnienia mają miarę zero.
Jednocześnie każdy singleton ma miarę zero.
Zatem nawet miara probabilistyczna może być atomowa,
chociaż żaden punkt nie ma dodatniej miary. Wniosek
``miara bez atomów wtedy i tylko wtedy, gdy wszystkie
singletony mają miarę zero'' jest bez dodatkowych
założeń fałszywy.
\end{przyklad}

\begin{twierdzenie}[Atomy miar $\sigma$-skończonych mają skończoną masę]
\label{thm:atom-sigma-finite}
Jeżeli $\mu$ jest miarą $\sigma$-skończoną na
$(X,\mathcal F)$ i $A\in\mathcal F$ jest atomem $\mu$,
to $0<\mu(A)<\infty$.
\end{twierdzenie}
\begin{proof}
Niech $X=\bigcup_nE_n$ z $\mu(E_n)<\infty$.
Gdyby wszystkie $A\cap E_n$ miały miarę zero,
to $\mu(A)=0$ na mocy przeliczalnej podaddytywności.
Istnieje więc $n$ z $\mu(A\cap E_n)>0$.
Atomowość $A$ implikuje $\mu(A\setminus E_n)=0$.
Otrzymujemy
$\mu(A)=\mu(A\cap E_n)+\mu(A\setminus E_n)
=\mu(A\cap E_n)\leq\mu(E_n)<\infty$.
\end{proof}

\begin{wniosek}[Przeliczalność rozłącznych atomów]
\label{cor:countably-many-disjoint-atoms}
Niech $\mu$ będzie miarą $\sigma$-skończoną na
$(X,\mathcal F)$, a $(A_i)_{i\in I}$ rodziną parami rozłącznych
atomów miary $\mu$. Wtedy zbiór indeksów $I$ jest co najwyżej
przeliczalny.
\end{wniosek}
\begin{proof}
Wybierzmy $X=\bigcup_{n\ge1}E_n$ z $\mu(E_n)<\infty$.
Z twierdzenia~\ref{thm:atom-sigma-finite} każdy $A_i$ ma
miarę dodatnią i skończoną. Istnieje więc $n$ takie, że
$\mu(A_i\cap E_n)>0$; z atomowości wynika wówczas
$\mu(A_i\cap E_n)=\mu(A_i)$. Dla ustalonych $n,k\ge1$
rodzina tych indeksów $i$, dla których
$\mu(A_i\cap E_n)\ge1/k$, jest skończona:
w przeciwnym przypadku suma miar dowolnie wielu
rozłącznych fragmentów przekroczyłaby $\mu(E_n)$.
Każdy indeks należy do przynajmniej jednej takiej
rodziny, więc $I$ jest przeliczalną sumą zbiorów
skończonych.
\end{proof}

\begin{uwaga}[Co daje $\sigma$-skończoność?]
$\sigma$-skończoność nie jest jedynie wygodnym
wzmocnieniem skończoności. Stanowi istotne założenie
jednoznaczności rozszerzenia miary, wielu wersji
twierdzenia Fubiniego oraz twierdzenia Radona--Nikodyma.
Nie należy jednak automatycznie zastępować
$\sigma$-skończoności półskończonością: są to różne hipotezy.
\end{uwaga}

\section{Synteza rozdziału i perspektywa dalszych konstrukcji}

Z rozdziałów \ref{ch:przestrzenie-mierzalne} i
\ref{ch:miary-wlasnosci} otrzymaliśmy następujący obraz.
$\sigma$-algebra kontroluje stabilność zbiorów mierzalnych
przy granicach przeliczalnych, zaś miara przekształca
rozłączne sumy w szeregi liczb nieujemnych. Te dwa
aksjomaty wystarczają do otrzymania monotoniczności,
podaddytywności, ciągłości z dołu i z góry oraz pierwszych
argumentów typu ``prawie wszędzie''.

Pozostaje jednak zasadnicze pytanie konstrukcyjne:
jak otrzymać miarę o zadanych wartościach na prostych
zbiorach, na przykład na prostokątach w $\mathbb R^d$?
Twierdzenia o jednoznaczności mówią, że pewne dane
wyznaczają miarę co najwyżej jedną; nie zapewniają
jej istnienia. W następnym rozdziale przejdziemy
do miar zewnętrznych i konstrukcji Carathéodory'ego,
która pozwoli przejść od prostych danych do prawdziwej
przeliczalnie addytywnej miary.
